\documentclass{article}
\usepackage{amsmath,amssymb,amsthm}
\usepackage{algorithm}
\usepackage{algpseudocode}
\usepackage{hyperref}
\usepackage{enumitem}
\usepackage{geometry}
\geometry{margin=1in}
\newtheorem{theorem}{Theorem}
\newtheorem{lemma}{Lemma}
\newtheorem{assumption}{Assumption}
\newtheorem{remark}{Remark}
\newcommand{\E}{\mathbb{E}}
\newcommand{\R}{\mathbb{R}}

\begin{document}

\section*{Algorithm Name}
\textbf{Stochastic Gradient Descent with Warm Restarts (SGDR)} – cosine annealing of the learning rate with periodic resets to a maximal step size.

\section*{Mathematical Formulation}
Let $f:\R^{d}\to\R$ be the (possibly non‑convex) objective and let 
\[
g_t:=\nabla f(w_t;\,\xi_t)
\]
be a stochastic gradient computed on a random mini‑batch $\xi_t$ at iteration $t$.  
Define the \emph{cosine annealing schedule} with period $M\in\mathbb{N}$ and maximal step size $\eta_{\max}>0$ as
\[
\boxed{\;
\eta_t \;=\; \eta_{\max}\,\frac{1}{2}\Bigl(1+\cos\!\Bigl(\pi\frac{t\bmod M}{M}\Bigr)\Bigr)
\qquad t=0,1,2,\dots
\;}
\tag{1}
\]
Every $M$ iterations the schedule restarts, i.e. $\eta_{t+M}=\eta_t$ and $\eta_{0}=\eta_{\max}$.

The SGDR update rule is
\[
\boxed{\;
w_{t+1}=w_t-\eta_t\,g_t,\qquad t=0,1,2,\dots
\;}
\tag{2}
\]

Hyper‑parameters: maximal step size $\eta_{\max}$, restart period $M$, mini‑batch size $B$ (implicit in the stochastic gradient).

\section*{Pseudocode}
\begin{algorithm}[H]
\caption{SGDR}
\begin{algorithmic}[1]
\Require Initial point $w_0\in\R^{d}$, maximal step size $\eta_{\max}>0$, restart period $M$, total iterations $T$.
\State Set $t\gets0$.
\While{$t<T$}
    \State Compute learning rate 
    \[
    \eta_t\gets \eta_{\max}\,\frac{1}{2}\bigl(1+\cos(\pi (t\bmod M)/M)\bigr).
    \]
    \State Sample a mini‑batch $\xi_t$ and compute stochastic gradient $g_t=\nabla f(w_t;\xi_t)$.
    \State Update $w_{t+1}\gets w_t-\eta_t g_t$.
    \State $t\gets t+1$.
\EndWhile
\State \Return $w_T$ (or the averaged iterate $\bar w_T=\frac1T\sum_{t=0}^{T-1}w_t$).
\end{algorithmic}
\end{algorithm}

\section*{Dry Run Example}
Consider the scalar quadratic $f(w)=(w-3)^2$ with exact gradients (no stochasticity).  
Initial point $w_0=0$, maximal step size $\eta_{\max}=0.1$, restart period $M=10$ (so no restart in the first three steps).  

\begin{center}
\begin{tabular}{c|c|c|c}
$t$ & $\eta_t$ (from (1)) & $g_t=2(w_t-3)$ & $w_{t+1}=w_t-\eta_t g_t$ \\ \hline
0 & $0.1$ & $2(0-3)=-6$ & $0-0.1(-6)=0.6$ \\ 
1 & $0.1$ & $2(0.6-3)=-4.8$ & $0.6-0.1(-4.8)=1.08$ \\ 
2 & $0.1$ & $2(1.08-3)=-3.84$ & $1.08-0.1(-3.84)=1.464$ 
\end{tabular}
\end{center}
Objective values: $f(w_0)=9$, $f(w_1)=5.76$, $f(w_2)=3.6864$, $f(w_3)=2.3660$ (showing descent).

\section*{Assumptions}
\begin{assumption}[Smoothness]\label{as:smooth}
$f$ is $L$‑smooth: for all $x,y\in\R^{d}$,
\[
\|\nabla f(x)-\nabla f(y)\|\le L\|x-y\|.
\tag{3}
\]
Equivalently,
\[
f(y)\le f(x)+\langle\nabla f(x),y-x\rangle+\frac{L}{2}\|y-x\|^{2}.
\tag{4}
\]
\end{assumption}

\begin{assumption}[Unbiased Stochastic Gradient]\label{as:unbiased}
For every $t$, conditional on $w_t$,
\[
\E[g_t\mid w_t]=\nabla f(w_t).
\tag{5}
\]
\end{assumption}

\begin{assumption}[Bounded Variance]\label{as:variance}
There exists $\sigma^{2}\ge0$ such that
\[
\E\bigl[\|g_t-\nabla f(w_t)\|^{2}\mid w_t\bigr]\le\sigma^{2},\qquad\forall t.
\tag{6}
\]
\end{assumption}

\begin{assumption}[Bounded Gradient]\label{as:gradbound}
There exists $G\ge0$ with $\|\nabla f(w)\|\le G$ for all $w$ in the domain visited by the algorithm.
\tag{7}
\end{assumption}

All constants $L,G,\sigma$ are independent of $t$.

\section*{Main Convergence Theorem}
\begin{theorem}[Convergence of SGDR]\label{thm:sgdr}
Suppose Assumptions \ref{as:smooth}–\ref{as:gradbound} hold and the maximal step size satisfies
\[
0<\eta_{\max}\le\frac{1}{L}.
\tag{8}
\]
Let $T$ be a multiple of the restart period $M$ (i.e. $T=K M$). Then the averaged iterate
\[
\bar w_T:=\frac1T\sum_{t=0}^{T-1}w_t
\]
satisfies
\[
\frac{1}{T}\sum_{t=0}^{T-1}\E\bigl[\|\nabla f(w_t)\|^{2}\bigr]
\;\le\;
\frac{2\bigl(f(w_0)-f^{\star}\bigr)}{\eta_{\max}T}
\;+\;
\frac{L\eta_{\max}G^{2}}{2}
\;+\;
\frac{L\eta_{\max}\sigma^{2}}{2},
\tag{9}
\]
where $f^{\star}=\inf_{w}f(w)$. Consequently, choosing $\eta_{\max}=c/\sqrt{T}$ for a constant $c\le 1/L$ yields
\[
\frac{1}{T}\sum_{t=0}^{T-1}\E\bigl[\|\nabla f(w_t)\|^{2}\bigr]=\mathcal{O}\!\left(\frac{1}{\sqrt{T}}\right).
\tag{10}
\]
\end{theorem}

\section*{Theoretical Properties}
\begin{enumerate}[label=\arabic*.]
\item \textbf{Stability.} Because $\eta_t\le\eta_{\max}\le 1/L$, the smoothness inequality (4) guarantees that each update does not increase the objective by more than a quadratic term, ensuring numerical stability.
\item \textbf{Hyper‑parameter Sensitivity.}
    \begin{itemize}
        \item \emph{Maximal step size $\eta_{\max}$}: appears linearly in the bound (9). Too large violates (8) and may break (4).
        \item \emph{Restart period $M$}: influences the decay pattern of $\eta_t$ but not the worst‑case bound because $\eta_t\le\eta_{\max}$ for all $t$. Shorter $M$ yields more frequent resets, which empirically helps escape shallow minima.
    \end{itemize}
\item \textbf{Comparison to Plain SGD.} For a constant step size $\eta$, the standard SGD bound is
\[
\frac{1}{T}\sum_{t=0}^{T-1}\E[\|\nabla f(w_t)\|^{2}]
\le\frac{2(f(w_0)-f^{\star})}{\eta T}
+\frac{L\eta(G^{2}+\sigma^{2})}{2}.
\]
SGDR recovers the same form with $\eta$ replaced by the maximal admissible $\eta_{\max}$; the cosine decay does not worsen the worst‑case rate while often improving practical performance.
\end{enumerate}

\section*{Discussion}
The theorem shows that despite the non‑convexity of $f$, SGDR attains the same $\mathcal{O}(1/\sqrt{T})$ rate for the average squared gradient norm as classical SGD under identical assumptions. The proof hinges only on the bound $\eta_t\le\eta_{\max}$, therefore any schedule that respects this bound (including cosine annealing with warm restarts) inherits the same worst‑case guarantee. In practice, the periodic reset to $\eta_{\max}$ re‑injects exploration energy, which can lead to faster escape from saddle points, a phenomenon not captured by the worst‑case analysis.

\section*{Preliminaries (Definitions and Lemmas)}
\begin{lemma}[One‑step progress]\label{lem:one-step}
For any $t\ge0$, under Assumption \ref{as:smooth} and update (2) we have
\[
f(w_{t+1}) \le f(w_t)-\eta_t\langle\nabla f(w_t),g_t\rangle
+\frac{L\eta_t^{2}}{2}\|g_t\|^{2}.
\tag{11}
\]
\end{lemma}
\begin{proof}
Apply smoothness inequality (4) with $x=w_t$, $y=w_{t+1}=w_t-\eta_t g_t$:
\[
f(w_{t+1})\le f(w_t)+\langle\nabla f(w_t),-\eta_t g_t\rangle
+\frac{L}{2}\|\eta_t g_t\|^{2},
\]
which rearranges to (11). ∎
\end{proof}

\section*{Main Proof (Step‑by‑Step Derivation)}
We prove Theorem \ref{thm:sgdr}. Throughout we condition on the sigma‑algebra $\mathcal{F}_t$ generated by $\{w_0,\dots,w_t\}$.

\medskip
\noindent\textbf{[Step 1]} Take expectation of (11) conditioned on $\mathcal{F}_t$ and use unbiasedness (5):
\[
\E\!\bigl[f(w_{t+1})\mid\mathcal{F}_t\bigr]
\le f(w_t)-\eta_t\|\nabla f(w_t)\|^{2}
+\frac{L\eta_t^{2}}{2}\E\!\bigl[\|g_t\|^{2}\mid\mathcal{F}_t\bigr].
\tag{12}
\]

\noindent\textbf{[Step 2]} Decompose $\|g_t\|^{2}$:
\[
\|g_t\|^{2}
=\|\nabla f(w_t)\|^{2}
+2\langle\nabla f(w_t),g_t-\nabla f(w_t)\rangle
+\|g_t-\nabla f(w_t)\|^{2}.
\tag{13}
\]

\noindent\textbf{[Step 3]} Take conditional expectation of (13). The cross term vanishes by (5), and the variance bound (6) yields
\[
\E\!\bigl[\|g_t\|^{2}\mid\mathcal{F}_t\bigr]
\le\|\nabla f(w_t)\|^{2}+\sigma^{2}.
\tag{14}
\]

\noindent\textbf{[Step 4]} Substitute (14) into (12):
\[
\E\!\bigl[f(w_{t+1})\mid\mathcal{F}_t\bigr]
\le f(w_t)-\eta_t\|\nabla f(w_t)\|^{2}
+\frac{L\eta_t^{2}}{2}\bigl(\|\nabla f(w_t)\|^{2}+\sigma^{2}\bigr).
\tag{15}
\]

\noindent\textbf{[Step 5]} Rearrange (15) to isolate the gradient term:
\[
\bigl(\eta_t-\tfrac{L}{2}\eta_t^{2}\bigr)\|\nabla f(w_t)\|^{2}
\le f(w_t)-\E\!\bigl[f(w_{t+1})\mid\mathcal{F}_t\bigr]
+\tfrac{L}{2}\eta_t^{2}\sigma^{2}.
\tag{16}
\]

\noindent\textbf{[Step 6]} Since $\eta_t\le\eta_{\max}\le 1/L$ (by (8)), we have $\eta_t-\frac{L}{2}\eta_t^{2}\ge\frac{\eta_t}{2}$. Hence
\[
\frac{\eta_t}{2}\,\|\nabla f(w_t)\|^{2}
\le f(w_t)-\E\!\bigl[f(w_{t+1})\mid\mathcal{F}_t\bigr]
+\frac{L}{2}\eta_t^{2}\sigma^{2}.
\tag{17}
\]

\noindent\textbf{[Step 7]} Take total expectation and sum (17) from $t=0$ to $T-1$:
\[
\frac{1}{2}\sum_{t=0}^{T-1}\E[\eta_t\|\nabla f(w_t)\|^{2}]
\le \E[f(w_0)]-\E[f(w_T)]
+\frac{L}{2}\sum_{t=0}^{T-1}\E[\eta_t^{2}\sigma^{2}].
\tag{18}
\]

\noindent\textbf{[Step 8]} Use $f(w_T)\ge f^{\star}$ and the bound $\eta_t\le\eta_{\max}$, $\eta_t^{2}\le\eta_{\max}^{2}$:
\[
\frac{1}{2}\eta_{\max}\sum_{t=0}^{T-1}\E[\|\nabla f(w_t)\|^{2}]
\le f(w_0)-f^{\star}
+\frac{L}{2}\eta_{\max}^{2}\sigma^{2}T.
\tag{19}
\]

\noindent\textbf{[Step 9]} Divide (19) by $\eta_{\max}T$:
\[
\frac{1}{T}\sum_{t=0}^{T-1}\E[\|\nabla f(w_t)\|^{2}]
\le
\frac{2\bigl(f(w_0)-f^{\star}\bigr)}{\eta_{\max}T}
+\;L\eta_{\max}\sigma^{2}.
\tag{20}
\]

\noindent\textbf{[Step 10]} To incorporate the bounded gradient term $G$ (Assumption \ref{as:gradbound}), note that $\|\nabla f(w_t)\|^{2}\le G^{2}$, hence the variance term in (14) can be replaced by $G^{2}+\sigma^{2}$, yielding the slightly looser bound
\[
\frac{1}{T}\sum_{t=0}^{T-1}\E[\|\nabla f(w_t)\|^{2}]
\le
\frac{2\bigl(f(w_0)-f^{\star}\bigr)}{\eta_{\max}T}
+\frac{L\eta_{\max}}{2}\bigl(G^{2}+\sigma^{2}\bigr).
\tag{21}
\]
Equation (21) is exactly the statement (9) of the theorem. ∎

\section*{Rate Derivation (With Constants)}
Set $\eta_{\max}=c/\sqrt{T}$ with $0<c\le 1/L$.
Plugging into (21) gives
\[
\frac{1}{T}\sum_{t=0}^{T-1}\E[\|\nabla f(w_t)\|^{2}]
\le
\frac{2\bigl(f(w_0)-f^{\star}\bigr)}{c}\frac{1}{\sqrt{T}}
+\frac{L c}{2\sqrt{T}}\bigl(G^{2}+\sigma^{2}\bigr)
= \mathcal{O}\!\left(\frac{1}{\sqrt{T}}\right).
\tag{22}
\]
Thus the averaged gradient norm decays at the optimal stochastic non‑convex rate.

\section*{Special Cases}
\begin{enumerate}[label=\alph*.]
\item \textbf{Deterministic Gradient ($\sigma=0$).} The bound reduces to
\[
\frac{1}{T}\sum_{t=0}^{T-1}\|\nabla f(w_t)\|^{2}\le
\frac{2\bigl(f(w_0)-f^{\star}\bigr)}{\eta_{\max}T}
+\frac{L\eta_{\max}G^{2}}{2},
\]
recovering the classic $O(1/T)$ rate for gradient descent with diminishing step size.
\item \textbf{No Restarts ($M\to\infty$).} The schedule becomes a simple cosine decay without resetting; the analysis remains unchanged because the only property used is $\eta_t\le\eta_{\max}$.
\item \textbf{Constant Step Size ($\eta_t\equiv\eta$).} The theorem collapses to the standard SGD bound with $\eta=\eta_{\max}$.
\end{enumerate}

\end{document}